\documentclass[11pt,a4paper]{article}
\usepackage[margin=22mm,headheight=15pt]{geometry}
\usepackage{fontspec}
\usepackage{polyglossia}
\setdefaultlanguage{russian}
\setotherlanguage{english}
\setmainfont{Arial}
\setsansfont{Arial}
\setmonofont{Menlo}
\newfontfamily\cyrillicfont{Arial}
\newfontfamily\cyrillicfonttt{Menlo}
\usepackage{unicode-math}
\setmathfont{STIX Two Math}
\usepackage{xcolor,booktabs,array,enumitem,graphicx,amsmath}
\usepackage[most]{tcolorbox}
\usepackage{fancyhdr}
\usepackage[colorlinks=true,linkcolor=blue,urlcolor=blue]{hyperref}
\definecolor{navy}{HTML}{102044}
\definecolor{accent}{HTML}{1168E8}
\definecolor{pale}{HTML}{F3F8FF}
\definecolor{yellow}{HTML}{FFDF55}
\definecolor{muted}{HTML}{4D627F}
\setlength{\parindent}{0pt}
\setlength{\parskip}{0.55em}
\setlist[enumerate]{leftmargin=*,itemsep=0.35em,topsep=0.4em}
\setlist[itemize]{leftmargin=*,itemsep=0.25em,topsep=0.3em}
\pagestyle{fancy}
\fancyhf{}
\fancyhead[L]{\small\color{navy} Введение в машинное обучение · Пара 2}
\fancyhead[R]{\small\color{accent} Методичка}
\fancyfoot[C]{\small\thepage}
\renewcommand{\headrulewidth}{0.4pt}
\newtcolorbox{saybox}{colback=pale,colframe=accent,boxrule=0pt,leftrule=3pt,arc=1mm,breakable,left=3mm,right=3mm,top=2mm,bottom=2mm}
\newtcolorbox{warnbox}{colback=yellow!18,colframe=yellow!80!black,boxrule=0pt,leftrule=3pt,arc=1mm,breakable,left=3mm,right=3mm,top=2mm,bottom=2mm}
\newcommand{\fig}[2]{%
  \begin{center}
  \includegraphics[width=0.88\linewidth]{#1}\\[1mm]
  {\small\color{muted}#2\par}
  \end{center}}
\newcounter{deflist}
\newcommand{\Def}[1]{\refstepcounter{deflist}\par\noindent{\color{navy}\textbf{Определение \thedeflist.}}\ #1\par}
\newcounter{stmtlist}
\newcommand{\Stmt}[1]{\refstepcounter{stmtlist}\par\noindent{\color{navy}\textbf{Утверждение \thestmtlist.}}\ #1\par}
\newcommand{\Ex}{\par\noindent{\color{navy}\textbf{Пример.}}\ }
\newcommand{\Rem}{\par\noindent{\color{navy}\textbf{Замечание.}}\ }
\begin{document}

{\color{accent}\bfseries ДГТУ · Введение в машинное обучение\par}
\vspace{4mm}
{\color{navy}\bfseries\fontsize{26}{32}\selectfont Линейная регрессия\par}
\vspace{3mm}
{\color{muted}\large Конспект лекции 2. Постановка задачи, метод наименьших квадратов, метрики качества и протокол оценивания.\par}
\vspace{2mm}

Изложение следует контуру курса линейных моделей: алгоритм, линейный по параметрам, квадратичный функционал, нормальные уравнения и ортогональность остатка. Градиентный спуск, регуляризация и логистическая регрессия на этой лекции не выводятся. Они названы лишь там, где иначе легко смешать разные задачи.

Сквозной числовой пример один: признак принимает значения \(1,2,3\), ответы равны \(1,2,2\). Все дроби ниже получены из него прямым счётом.

\section{Постановка задачи и обозначения}

Пусть имеется выборка
\[
(x_i,y_i),\qquad i=1,\ldots,n,
\]
где \(y_i\in\mathbb R\) --- целевая переменная, а признаковое описание объекта без свободного члена есть \((x_{i1},\ldots,x_{id})\).

В матрицу плана \(X\) слева приписывают столбец единиц, если в модели есть свободный член. Тогда строка объекта имеет вид
\[
x_i^{\top}=(1,\ x_{i1},\ \ldots,\ x_{id})\in\mathbb R^{p},\qquad p=d+1,
\]
и \(X\in\mathbb R^{n\times p}\). Вектор параметров обозначается \(w\in\mathbb R^{p}\). Алгоритм
\[
a(x)=x^{\top}w
\]
сопоставляет объекту одно число. Это же число, вычисленное на \(i\)-м объекте обучающей выборки, записывают как \(\hat y_i\). Остаток
\[
r_i=y_i-\hat y_i
\]
есть часть ответа, которую алгоритм не воспроизвёл. Вектор всех ответов --- \(y\in\mathbb R^{n}\), вектор прогнозов --- \(\hat y=Xw\).

\Def{Задачей восстановления регрессии называется построение алгоритма \(a\), значение которого приближает целевую переменную. При фиксированном \(x\) величина \(y\) остаётся случайной, поэтому алгоритм предсказывает выбранную характеристику условного распределения \(y\mid x\), а не единственно возможный ответ.}

На этой лекции такой характеристикой служит условное математическое ожидание, и оно ищется в классе линейных по \(w\) функций:
\begin{equation}
\label{eq:cond-mean}
\mathbb E[y\mid x]=x^{\top}w.
\end{equation}

Классификацией называется другая задача: ответ принадлежит конечному множеству меток. Общность названия «регрессия» у логистической модели эту границу не снимает.

\section{Происхождение термина}

Ф.~Гальтон в 1885--1886 годах сопоставил рост родителей и рост взрослых сыновей. Рост матери предварительно умножался примерно на \(1{,}08\), чтобы привести его к мужской шкале; затем бралось среднее двух родителей. На график наносился средний рост сыновей в группе с данным ростом родителей. Эти средние легли близко к прямой с наклоном около \(2/3\).

\fig{figures/galton-line.png}{Мелкая точка --- отдельная семья, то есть один объект выборки. Крупная точка --- средний рост сыновей в группе с данным ростом родителей. Прямая этих условных средних имеет меньший наклон, чем биссектриса.}

Содержательно наклон \(2/3\) означает следующее. Отклонение родительского роста от среднего популяции переходит сыну не полностью, а примерно на две трети. Гальтон назвал это явление возвратом к середине, \emph{regression towards mediocrity}. Отсюда термин «регрессия».

При близких дисперсиях признака и ответа условное среднее записывается как
\begin{equation}
\label{eq:galton}
\mathbb E[y\mid x]=\mu+\beta(x-\mu),
\end{equation}
где \(\mu\) --- общее среднее данной записи, а \(\beta\) совпадает с коэффициентом корреляции. Пока связь не вырождена в точное равенство, \(|\beta|<1\), и предсказанное отклонение короче отклонения \(x\). Биссектриса соответствовала бы случаю \(\beta=1\): отклонение родителей переходит сыну целиком.

\Rem Если дисперсии различны, коэффициент в~\eqref{eq:galton} равен \(\beta=\rho\,\sigma_y/\sigma_x\) и может превосходить единицу. Возврат отклонения к середине в обычном смысле слова имеет место при \(0<\beta<1\).

Тот же механизм наблюдается вне биологии. Группа, отобранная по крайне высокому результату первого измерения, при повторном измерении в среднем ближе к центру распределения, если связь двух измерений неидеальна. Сдвиг группы является свойством условного среднего. Приписывать его произошедшему между измерениями вмешательству нельзя.

\begin{warnbox}
Метод наименьших квадратов к опытам Гальтона уже существовал: публикация А.-М.~Лежандра относится к 1805 году, К.~Ф.~Гаусс применял метод ранее. К.~Пирсон впоследствии ввёл коэффициент корреляции как формальный объект. Вывод метода наименьших квадратов Пирсону не принадлежит. Термин \emph{regression} закреплён за картиной условных средних у Гальтона.
\end{warnbox}

\section{Семейство моделей}

Слово «регрессия» обозначает несколько различных моделей. Ниже фиксируются те, которые нужны для дальнейших лекций. На этой лекции выводятся две первые.

\Def{Простой линейной регрессией называется модель \(a(x)=w_0+w_1 x\) с одним вещественным признаком.}

\Def{Множественной линейной регрессией называется модель
\[
a(x)=w_0+w_1 x_1+\cdots+w_d x_d
\]
с несколькими признаками. Значение алгоритма по-прежнему одно число.}

\Ex Пусть цена квартиры задаётся свободным членом \(20\) и весом площади \(2\). При площади \(40\)~м\textsuperscript{2}
\[
a=20+2\cdot 40=100.
\]
Если добавлен вес этажа \(3\), та же квартира на \(5\)-м этаже получает прогноз
\[
a=20+2\cdot 40+3\cdot 5=115.
\]
Изменилось число признаков. Размерность ответа осталась равной единице.

\Def{Полиномиальной регрессией по одному признаку называется линейная по параметрам модель, среди столбцов которой присутствуют степени этого признака. Для степени \(2\)
\[
a(x)=w_0+w_1 x+w_2 x^{2}.
\]
}

Модель остаётся линейной по \(w\): столбец \(x^{2}\) вычисляется до оценивания параметров и входит в скалярное произведение наравне с остальными. Нелинейность относится к графику зависимости от исходного признака.

\Ex Запись \(20+2x+0{,}01\,x^{2}\) приведена как иллюстрация вида модели, а не как результат оценивания. При \(x\in\{20,40,60\}\) она даёт значения \(64\), \(116\) и \(176\). Приращение признака одинаково и равно \(20\), приращения прогноза равны \(52\) и \(60\). График не является прямой. Оценивание таких столбцов выполняется тем же методом наименьших квадратов и относится к последующей лекции.

\Rem Логистическая регрессия оценивает вероятность метки класса. К задаче~\eqref{eq:cond-mean} она не сводится. Общим является слово в названии.

\begin{saybox}
Предметом лекции являются простая и множественная линейная регрессия. В обеих значение алгоритма есть одно число, и восстанавливается условное среднее.
\end{saybox}

\section{Среднее, медиана и квантиль}

Прежде чем строить зависимость от \(x\), полезно различить характеристики одномерной выборки. Пусть ответы упорядочены: \(y_{(1)}\le\cdots\le y_{(n)}\).

\Def{Выборочным средним называется \(\bar y=n^{-1}\sum_{i=1}^{n} y_i\).}

\Def{Выборочной медианой называется \(y_{((n+1)/2)}\) при нечётном \(n\) и
\[
\tfrac12\bigl(y_{(n/2)}+y_{(n/2+1)}\bigr)
\]
при чётном \(n\).}

\Def{Выборочным квантилем уровня \(\tau\in(0,1)\) в этом конспекте называется порядковая статистика \(y_{(k)}\), где \(k=\lceil \tau n\rceil\). Это наименьшее значение, для которого доля наблюдений, не превосходящих его, не меньше \(\tau\).}

\Ex Для выборки \(2,2,3,10\) вариационный ряд совпадает с исходной записью, \(n=4\).
\begin{align*}
\bar y&=\frac{2+2+3+10}{4}=\frac{17}{4}=4{,}25,\\[1mm]
\text{медиана}&=\frac{2+3}{2}=2{,}5,\\[1mm]
k=\lceil 0{,}9\cdot 4\rceil&=\lceil 3{,}6\rceil=4,\qquad y_{(4)}=10.
\end{align*}
Три характеристики различны. Среднее учитывает величину крайнего наблюдения: без значения \(10\) среднее оставшихся трёх равно \(7/3\). Медиана при замене \(10\) на \(100\) не изменится. Квантиль уровня \(0{,}9\) обязан накрыть не менее \(90\%\) наблюдений; три наблюдения из четырёх составляют лишь \(75\%\), поэтому порог совпадает с максимумом.

\fig{figures/numberline.png}{Три характеристики одной выборки: медиана \(2{,}5\), среднее \(4{,}25\), квантиль уровня \(0{,}9\), равный \(10\).}

\Rem При чётном \(n\) медиана как полусумма двух центральных порядковых статистик не обязана совпадать с квантилем уровня \(1/2\) в смысле определения выше. Для той же выборки \(k=\lceil 0{,}5\cdot 4\rceil=2\) и \(y_{(2)}=2\), тогда как медиана равна \(2{,}5\). Далее медиана понимается как полусумма, квантиль --- как порядковая статистика с номером \(\lceil\tau n\rceil\).

\Stmt{Константа, минимизирующая сумму квадратов отклонений, равна выборочному среднему. Константа, минимизирующая сумму модулей, равна выборочной медиане. Константа, минимизирующая сумму pinball-потерь уровня \(\tau\), равна квантилю этого уровня.}

Доказательство для квадрата будет следствием нормальных уравнений в модели из одного столбца единиц. Для модуля и pinball достаточно баланса односторонних приращений; оно приводится при соответствующих метриках. На этой лекции линейная регрессия восстанавливает характеристику из первой части утверждения, то есть среднее.

\section{Линейная модель и матрица плана}

Прогноз на одном объекте есть
\begin{equation}
\label{eq:model}
\hat y=w_0+w_1 x_1+\cdots+w_d x_d=x^{\top}w.
\end{equation}
Коэффициент \(w_j\) равен изменению прогноза при увеличении \(j\)-го признака на единицу и фиксированных остальных столбцах модели. Свободный член \(w_0\) равен прогнозу при нулевых значениях признаков, вошедших в модель.

Для выборки из \(n\) объектов те же параметры задают систему
\begin{equation}
\label{eq:matrix}
\hat y=Xw.
\end{equation}
Строка \(X\) соответствует объекту, столбец --- признаку. Первый столбец единиц делает свободный член обыкновенным параметром.

\Ex Сквозной пример в матричной записи:
\[
X=\begin{pmatrix}1&1\\1&2\\1&3\end{pmatrix},\qquad
y=\begin{pmatrix}1\\2\\2\end{pmatrix},\qquad
w=\begin{pmatrix}w_0\\w_1\end{pmatrix}.
\]
Первая строка читается как прогноз \(w_0+w_1\cdot 1\).

\Rem Отсутствие столбца единиц есть содержательное ограничение: алгоритм обязан обращаться в нуль при нулевых признаках. Для цены квартиры это означало бы нулевую цену при нулевой площади. Ограничение вводят только тогда, когда оно следует из постановки. В стандартной реализации свободный член включён по умолчанию.

Линейность модели есть линейность по параметрам. Замена исходных признаков на их степени, логарифмы или индикаторы категорий, выполненная до оценивания \(w\), класс не покидает: функционал по \(w\) остаётся квадратичным и выпуклым. Запись \(w_0+e^{w_1}x\) в этот класс не входит, поскольку параметр стоит внутри нелинейной функции.

\section{Категориальные признаки}

Вещественный код номинального признака, например номера района \(1,2,3\), не является допустимым столбцом. Линейная модель интерпретирует такой столбец как упорядоченный признак с равными шагами между уровнями.

\Def{One-hot-кодированием \(K\) уровней называется набор из \(K\) индикаторов: в строке объекта стоит одна единица, в столбце его уровня, и нули в остальных.}

\Def{One-cold-кодированием называется дополнение one-hot до единицы в каждой координате: нули и единицы меняются местами. В строке остаётся ровно один нуль.}

Для трёх районов коды таковы.

\begin{center}
\begin{tabular}{@{}lcc@{}}
\toprule
Уровень & one-hot & one-cold\\
\midrule
Центр & \(1\ 0\ 0\) & \(0\ 1\ 1\)\\
Север & \(0\ 1\ 0\) & \(1\ 0\ 1\)\\
Юг & \(0\ 0\ 1\) & \(1\ 1\ 0\)\\
\bottomrule
\end{tabular}
\end{center}

\fig{figures/onehot.png}{Те же три уровня. Закрашенная клетка соответствует единице. One-cold получен из one-hot поклеточной заменой значения на противоположное.}

\Stmt{Полный one-hot и полный one-cold вместе со столбцом единиц линейно зависимы. В первом случае сумма индикаторов равна столбцу единиц. Во втором при \(K=3\) сумма столбцов равна удвоенному столбцу единиц, поскольку каждая строка содержит две единицы. Матрица \(X^{\top}X\) вырождена, и параметр \(w\) не единствен.}

Прогнозы \(Xw\) при этом могут оставаться единственными: к решению прибавляется вектор из ядра \(X\), и произведение \(Xw\) не меняется.

Рабочая спецификация со свободным членом состоит в удалении одного уровня. Удалённый уровень называется базовым. Коэффициент сохранённого уровня равен разности прогноза относительно базового уровня при фиксированных остальных столбцах.

\Ex Пусть базовым уровнем выбран Юг, свободный член равен \(100\), коэффициент Центра равен \(+20\), коэффициент Севера равен \(-10\). Тогда прогноз Юга равен \(100\), прогноз Центра равен \(120\), прогноз Севера равен \(90\). Число \(+20\) есть контраст «Центр против Юга», а не абсолютная цена Центра.

Словарь уровней оценивается только по обучающей выборке. Уровень, не встретившийся в ней, отдельного столбца не получает.

\section{Функционал наименьших квадратов}

\Def{Методом наименьших квадратов называется минимизация по \(w\) суммы квадратов остатков на той выборке, по которой оцениваются параметры:
\begin{equation}
\label{eq:rss}
\mathrm{RSS}(w)=\lVert y-Xw\rVert^{2}
=\sum_{i=1}^{n}(y_i-x_i^{\top}w)^{2}.
\end{equation}
Эмпирический риск \(Q(w)=n^{-1}\mathrm{RSS}(w)\) имеет тот же минимизатор: положительный множитель точку минимума не сдвигает.}

Квадрат остатка задаёт два свойства, которые следует разделять. Ошибка \(1\) входит в сумму как \(1\), ошибка \(3\) входит как \(9\): крупные отклонения доминируют. Кроме того, квадрат дифференцируем, и условие минимума записывается системой линейных уравнений.

\section{Вывод для одного признака}

Пусть \(a(x)=w_0+w_1 x\) и
\[
\mathrm{RSS}(w_0,w_1)=\sum_{i=1}^{n}(y_i-w_0-w_1 x_i)^{2}.
\]

Частная производная по свободному члену равна
\[
\frac{\partial\mathrm{RSS}}{\partial w_0}
=-2\sum_{i=1}^{n}(y_i-w_0-w_1 x_i).
\]
В критической точке сумма остатков обращается в нуль, откуда
\begin{equation}
\label{eq:intercept}
w_0=\bar y-w_1\bar x.
\end{equation}
Следовательно, прямая проходит через центр выборки \((\bar x,\bar y)\).

Частная производная по наклону равна
\[
\frac{\partial\mathrm{RSS}}{\partial w_1}
=-2\sum_{i=1}^{n}x_i(y_i-w_0-w_1 x_i).
\]
Подстановка~\eqref{eq:intercept} и тождество \(\sum(x_i-\bar x)=0\) дают
\begin{equation}
\label{eq:slope}
w_1=
\frac{\sum_{i=1}^{n}(x_i-\bar x)(y_i-\bar y)}
{\sum_{i=1}^{n}(x_i-\bar x)^{2}},
\qquad
w_0=\bar y-w_1\bar x.
\end{equation}
Числитель измеряет совместное отклонение признака и ответа от своих средних, знаменатель --- разброс признака. Общий множитель \(1/n\) или \(1/(n-1)\) в числителе и знаменателе сокращается и на оценку \(w_1\) не влияет.

\Rem В модели из одного столбца единиц наклона нет. Условие~\eqref{eq:intercept} превращается в \(w_0=\bar y\). Тем самым доказана первая часть утверждения о константе: для квадратичного функционала оптимальная константа есть выборочное среднее обучающей выборки.

Для суммы модулей производная в нуле не существует, и замкнутой формулы наклона, аналогичной~\eqref{eq:slope}, нет. Для константы сумма модулей убывает при сдвиге в сторону, где наблюдений больше. Баланс достигается, когда по обе стороны от константы лежит поровну наблюдений, то есть в медиане.

\section{Нормальные уравнения}

Для произвольного \(p\) градиент~\eqref{eq:rss} равен
\begin{equation}
\label{eq:grad}
\nabla\mathrm{RSS}(w)=-2X^{\top}(y-Xw).
\end{equation}
Каждая координата есть частная производная по своему параметру. Критическая точка задаётся системой
\begin{equation}
\label{eq:normal}
X^{\top}X\,w=X^{\top}y,
\end{equation}
называемой нормальными уравнениями. Для простой регрессии она раскрывается в
\[
\begin{cases}
n\,w_0+\bigl(\sum x_i\bigr)w_1=\sum y_i,\\[1mm]
\bigl(\sum x_i\bigr)w_0+\bigl(\sum x_i^{2}\bigr)w_1=\sum x_i y_i.
\end{cases}
\]
Первое уравнение есть условие нулевой суммы остатков. Второе есть ортогональность остатка столбцу признака.

Гессиан \(\mathrm{RSS}\) равен \(2X^{\top}X\). Для любого направления \(d\in\mathbb R^{p}\)
\[
d^{\top}(2X^{\top}X)\,d=2\lVert Xd\rVert^{2}\ge 0.
\]
Функционал выпуклый, седловых точек нет, критическая точка является глобальным минимумом. При полном столбцовом ранге минимум единствен и
\begin{equation}
\label{eq:inverse}
\hat w=(X^{\top}X)^{-1}X^{\top}y.
\end{equation}
Формула~\eqref{eq:inverse} предполагает обратимость \(X^{\top}X\). В вычислениях обратную матрицу не формируют: систему решают разложением QR или SVD. Число обусловленности \(X^{\top}X\) равно квадрату числа обусловленности \(X\), поэтому явное обращение усиливает погрешность правой части.

\Rem При неполном ранге минимизаторов много, однако вектор прогнозов \(X\hat w\) у них общий: это ортогональная проекция \(y\) на линейную оболочку столбцов \(X\). Прибавление к \(\hat w\) произвольного вектора из ядра \(X\) прогноз на обучающих объектах не изменяет.

\section{Сквозной пример}

Пусть \(x\in\{1,2,3\}\) и \(y\in\{1,2,2\}\). Тогда
\[
\bar x=\frac{1+2+3}{3}=2,\qquad
\bar y=\frac{1+2+2}{3}=\frac{5}{3}.
\]
Отклонения признака равны \(-1,0,1\), отклонения ответа равны \(-2/3,\ 1/3,\ 1/3\). По формуле~\eqref{eq:slope}
\[
w_1=\frac{(-1)(-2/3)+(1)(1/3)}{1+1}=\frac{1}{2},\qquad
w_0=\frac{5}{3}-\frac12\cdot 2=\frac{2}{3}.
\]

Те же значения даёт система. Здесь \(\sum x=6\), \(\sum x^{2}=14\), \(\sum y=5\), \(\sum xy=11\), поэтому
\[
\begin{cases}
3w_0+6w_1=5,\\
6w_0+14w_1=11.
\end{cases}
\]
Из первого уравнения \(w_0=5/3-2w_1\). Подстановка во второе:
\begin{align*}
6\bigl(5/3-2w_1\bigr)+14w_1&=11,\\
10+2w_1&=11,\\
w_1=1/2,&\qquad w_0=2/3.
\end{align*}

Прогнозы и остатки приведены в таблице.

\begin{center}
\begin{tabular}{@{}ccccc@{}}
\toprule
\(x\) & \(y\) & \(2/3+x/2\) & остаток & квадрат\\
\midrule
\(1\) & \(1\) & \(7/6\) & \(-1/6\) & \(1/36\)\\
\(2\) & \(2\) & \(5/3\) & \(+1/3\) & \(4/36\)\\
\(3\) & \(2\) & \(13/6\) & \(-1/6\) & \(1/36\)\\
\bottomrule
\end{tabular}
\end{center}

Сумма остатков равна нулю. Это следствие столбца единиц, а не частного подбора чисел. Сумма квадратов остатков и сумма квадратов отклонений ответа от среднего равны
\begin{align*}
\mathrm{RSS}&=\frac{1}{36}+\frac{4}{36}+\frac{1}{36}=\frac{1}{6},\\[1mm]
\mathrm{SST}&=\Bigl(-\frac23\Bigr)^{2}+\Bigl(\frac13\Bigr)^{2}+\Bigl(\frac13\Bigr)^{2}=\frac{2}{3}.
\end{align*}

\begin{saybox}
Вычисления выполнены на той же выборке, по которой получен \(\hat w\). Они описывают подогнанную прямую. Ошибку на новых объектах они не оценивают.
\end{saybox}

\section{Геометрическая интерпретация}

Равенство \(X^{\top}(y-X\hat w)=0\), эквивалентное~\eqref{eq:normal}, означает ортогональность остатка каждому столбцу \(X\). Вектор \(y\) рассматривается как точка пространства \(\mathbb R^{n}\). Столбцы \(X\) порождают подпространство размерности не выше \(p\). Метод наименьших квадратов выбирает в этом подпространстве точку, ближайшую к \(y\) в евклидовой норме. Ближайшая точка есть основание перпендикуляра.

\fig{figures/projection.png}{Остаток ортогонален линейной оболочке столбцов. Прямой угол относится к пространству ответов \(\mathbb R^{n}\), а не к углу прямой с осью признака на плоскости \((x,y)\).}

Для центрированного вектора ответов теорема Пифагора даёт разложение
\begin{equation}
\label{eq:pythagoras}
\lVert y-\bar y\,\mathbf 1\rVert^{2}
=\lVert X\hat w-\bar y\,\mathbf 1\rVert^{2}
+\lVert y-X\hat w\rVert^{2}.
\end{equation}
Левая часть есть \(\mathrm{SST}\). Второе слагаемое справа есть \(\mathrm{RSS}\). Первое слагаемое справа есть объяснённая сумма квадратов. На сквозном примере \(2/3=1/2+1/6\).

\Def{Коэффициентом детерминации на выборке, по которой он вычислен, называется доля объяснённой суммы квадратов:
\begin{equation}
\label{eq:r2}
R^{2}
=1-\frac{\mathrm{RSS}}{\mathrm{SST}}
=1-\frac{\sum_{i}(y_i-\hat y_i)^{2}}{\sum_{i}(y_i-\bar y)^{2}}.
\end{equation}
}

В сквозном примере \(R^{2}=1-(1/6)/(2/3)=3/4\). Величина равна доле суммы квадратов отклонений от выбранного среднего. Долей верно восстановленных ответов и вероятностью она не является. Значение \(0\) означает совпадение суммы квадратов с константой \(\bar y\). Значение \(1\) означает нулевые остатки. На объектах, не участвовавших в оценивании, \(R^{2}\) может быть отрицательным: алгоритм уступает выбранной константе.

\section{Единственность параметров и мультиколлинеарность}

Если два столбца \(X\) линейно зависимы, матрица \(X^{\top}X\) необратима. Множество параметров, доставляющих один и тот же прогноз \(Xw\), бесконечно. Интерпретация отдельного коэффициента как изменения прогноза «при прочих равных» в этом случае недоступна: прочие столбцы нельзя зафиксировать, не изменив зависимый столбец.

При близкой к вырожденной зависимости параметр формально единствен, но неустойчив. Если ошибки сферичны и ранг полон,
\[
\mathrm{Var}(\hat w)=\sigma^{2}(X^{\top}X)^{-1}.
\]
Малое собственное число \(X^{\top}X\) увеличивает дисперсию соответствующей линейной комбинации коэффициентов. Прогноз \(X\hat w\) при этом может оставаться устойчивым: неустойчива речь об отдельном весе, а не обязательно сама проекция.

\Ex Площадь и число комнат в одной спецификации, как правило, сильно зависимы. Коэффициент при площади не интерпретируют отдельно от коэффициента при числе комнат.

Численное обращение \(X^{\top}X\) усиливает погрешность входных данных. Стандартные реализации ищут решение наименьших квадратов через QR или SVD и среди минимизаторов одинаковой суммы квадратов выбирают вектор \(w\) минимальной евклидовой нормы. Стабилизация зависимых столбцов штрафом за норму \(w\) на этой лекции не рассматривается.

\section{Вероятностная модель ошибки}

Пусть \(y=Xw+\varepsilon\) и координаты \(\varepsilon\) независимы и распределены нормально с нулевым средним и дисперсией \(\sigma^{2}\). Плотность выборки пропорциональна
\[
\exp\bigl(-\lVert y-Xw\rVert^{2}/(2\sigma^{2})\bigr).
\]
Максимум правдоподобия по \(w\) совпадает с минимумом~\eqref{eq:rss}. Оценка дисперсии по максимуму правдоподобия равна \(\mathrm{RSS}/n\). Несмещённая оценка при полном ранге и \(p\) столбцах равна \(\mathrm{RSS}/(n-p)\). В сквозном примере вторая величина есть \((1/6)/(3-2)=1/6\). Выбор делителя влияет на оценку \(\sigma^{2}\) и не влияет на \(\hat w\).

Если координаты ошибки имеют распределение Лапласа, максимум правдоподобия совпадает с минимумом суммы модулей. Оптимальная константа тогда является медианой. Линейной системы вида~\eqref{eq:normal} эта задача не порождает.

\fig{figures/outlier-lines.png}{Сплошная прямая соответствует квадратичному функционалу и смещена к удалённому наблюдению. Пунктир соответствует сумме модулей. Рисунок схематичен и не является оценкой по конкретной выборке.}

Квадратичный штраф растёт как квадрат расстояния и потому сильно сдвигает прямую к удалённой точке. Модульный штраф растёт линейно; положение прямой определяется главным образом тем, по какую сторону остатка лежит большинство объектов.

\subsection{Теорема Гаусса--Маркова}

\Stmt{Предположим, что выполнены четыре условия.
\begin{enumerate}
\item Спецификация верна: \(\mathbb E[y\mid X]=Xw\).
\item Регрессоры экзогенны: \(\mathbb E[\varepsilon\mid X]=0\).
\item Ошибки сферичны: \(\mathrm{Var}(\varepsilon\mid X)=\sigma^{2} I\). Дисперсия постоянна, ошибки некоррелированы.
\item Столбцы \(X\) линейно независимы.
\end{enumerate}
Тогда оценка метода наименьших квадратов является лучшей линейной несмещённой оценкой параметра \(w\): среди оценок, линейных по \(y\) и несмещённых, она имеет наименьшую дисперсию. Доказательство на лекции не приводится.}

Нормальность \(\varepsilon\) в условия не входит. Она нужна для совпадения метода наименьших квадратов с максимумом правдоподобия. Сравнение в утверждении идёт только с линейными несмещёнными оценками, не со всеми мыслимыми оценками.

Нарушение сферичности при сохранении экзогенности часто оставляет \(\hat w\) несмещённой и лишает её свойства минимальной дисперсии в указанном классе. Обычная формула \(\sigma^{2}(X^{\top}X)^{-1}\) в этом случае дисперсию не описывает. Нарушение экзогенности, в частности пропуск фактора, входящего и в столбцы, и в ошибку, смещает \(\hat w\).

\begin{saybox}
Коэффициент \(w_j\) отвечает на вопрос об изменении прогноза при увеличении \(j\)-го столбца на единицу внутри выбранной спецификации. Решением системы~\eqref{eq:normal} причинная интерпретация не обеспечивается.
\end{saybox}

\section{Интерпретация коэффициентов}

Корректная формулировка такова: при фиксированных остальных столбцах модели увеличение \(j\)-го признака на единицу изменяет прогноз на \(w_j\).

Три уточнения обязательны.

Если столбец стандартизован, единица признака равна одному его стандартному отклонению. Коэффициент измеряет изменение прогноза на это отклонение. В исходных единицах, например в рублях на метр, он уже не выражен.

Коэффициент индикатора уровня равен контрасту с базовым уровнем, удалённым из матрицы. В примере с районами значение \(+20\) означало превышение прогноза Центра над прогнозом Юга.

При близкой зависимости столбцов отдельный коэффициент не отделяют от коэффициентов тех столбцов, с которыми он движется совместно.

\section{Функция потерь и метрика качества}

\Def{Функцией потерь называется величина, по которой на обучающей выборке ищется параметр \(w\). Метрикой качества называется величина, по которой уже выбранный алгоритм сравнивают на объектах, не участвовавших в этом поиске.}

Формулы могут совпадать. Роли от этого не совпадают. Оценивание по квадрату и отчёт по среднему модулю образуют допустимую пару, если протокол явно называет и функцию потерь, и метрику. Ошибки \(1\) и \(3\) входят в квадрат как \(1\) и \(9\), в модуль --- как \(1\) и \(3\).

Ниже метрики вычислены на сквозном примере, то есть на трёх обучающих объектах. Константа \(\bar y=5/3\) есть метод наименьших квадратов в модели из столбца единиц. По среднему квадрату прямая на этой выборке не может уступить константе: пространство столбцов расширено, функционал тот же. Для прочих метрик такого соотношения нет.

\subsection{Средний квадрат и его корень}

\[
\mathrm{MSE}=\frac1n\sum_{i=1}^{n}(y_i-\hat y_i)^{2},\qquad
\mathrm{RMSE}=\sqrt{\mathrm{MSE}}.
\]

\fig{figures/loss-mse.png}{Квадратичный штраф. Отклонение \(3\) даёт вклад \(9\) до усреднения.}

В сквозном примере \(\mathrm{RSS}=1/6\), поэтому у прямой \(\mathrm{MSE}=(1/6)/3=1/18\) и \(\mathrm{RMSE}=\sqrt{1/18}\). У константы \(5/3\) сумма квадратов равна \(2/3\), откуда \(\mathrm{MSE}=2/9\). Корень возвращает порядок единиц целевой переменной. Средним абсолютным отклонением \(\mathrm{RMSE}\) не является: крупный остаток по-прежнему проходит через квадрат.

\subsection{Средний модуль}

\[
\mathrm{MAE}=\frac1n\sum_{i=1}^{n}\lvert y_i-\hat y_i\rvert.
\]
Величина выражена в единицах целевой переменной. Для отклонений \(1\), \(2\) и \(3\) получается \((1+2+3)/3=2\).

\fig{figures/loss-mae.png}{Модульный штраф. Отклонение \(3\) даёт вклад \(3\).}

В сквозном примере модули остатков прямой равны \(1/6\), \(1/3\), \(1/6\). Их сумма есть \(2/3\), следовательно \(\mathrm{MAE}=2/9\). У константы \(5/3\) модули равны \(2/3\), \(1/3\), \(1/3\), сумма равна \(4/3\) и \(\mathrm{MAE}=4/9\). Выигрыш прямой по модулю в этом примере есть. Метод наименьших квадратов его не гарантирует, поскольку параметр искался по другому функционалу.

Оптимальность медианы для константы видна из одностороннего приращения. Сдвиг константы на малую величину \(\delta>0\) уменьшает вклад каждого наблюдения справа на \(\delta\) и увеличивает вклад каждого наблюдения слева на \(\delta\). Сумма убывает, пока левее константы наблюдений меньше, чем правее. Равенство сторон есть определение медианы при нечётном числе наблюдений; при чётном числе наблюдений минимумом является любая точка отрезка между двумя центральными порядковыми статистиками, и в качестве оценки берут его середину.

\subsection{Коэффициент детерминации и его поправка}

Определение дано формулой~\eqref{eq:r2}. В сквозном примере
\[
R^{2}=1-\frac{1/6}{2/3}=\frac34
\]
для прямой и \(R^{2}=0\) для константы \(5/3\).

\fig{figures/r2-bars.png}{Сумма квадратов вокруг среднего равна \(2/3\), вокруг прямой --- \(1/6\).}

\begin{warnbox}
Значение \(3/4\) есть доля суммы квадратов на трёх обучающих объектах. Процедура \texttt{r2\_score} подставляет в знаменатель среднее того вектора \(y\), который ей передан. Если передан контрольный вектор, среднее берётся по контролю. Сравнение с обучающим средним есть другая процедура. В протоколе фиксируют, какое среднее стоит в знаменателе.
\end{warnbox}

Скорректированный коэффициент
\begin{equation}
\label{eq:adj}
R^{2}_{\mathrm{adj}}=1-(1-R^{2})\frac{n-1}{n-p}
\end{equation}
вводит штраф за число столбцов \(p\), включая столбец единиц. В сквозном примере \(n=3\), \(p=2\):
\[
R^{2}_{\mathrm{adj}}=1-\Bigl(1-\frac34\Bigr)\frac{2}{1}=\frac12.
\]
Поправка вычисляется на той же выборке, по которой получен \(R^{2}\). Отложенную выборку она не заменяет. Спецификации с различными преобразованиями целевой переменной сравнивать по~\eqref{eq:adj} нельзя: знаменатель \(\mathrm{SST}\) живёт в разных шкалах.

\subsection{Относительные ошибки}

\[
\mathrm{MAPE}
=\frac1n\sum_{i=1}^{n}\frac{\lvert y_i-\hat y_i\rvert}{\lvert y_i\rvert}.
\]
Вклад абсолютного остатка обратно пропорционален модулю ответа. При \(y_i=0\) выражение не определено.

В сквозном примере для прямой
\[
\frac{1/6}{1}+\frac{1/3}{2}+\frac{1/6}{2}
=\frac{1}{6}+\frac{1}{6}+\frac{1}{12}=\frac{5}{12},
\qquad
\mathrm{MAPE}=\frac{5}{36}.
\]
Для константы \(5/3\) сумма трёх отношений равна \(1\), откуда \(\mathrm{MAPE}=1/3\).

На одном объекте прогноз «вдвое выше» и прогноз «вдвое ниже» дают одно и то же значение. На выборке с различными \(|y_i|\) оптимальная константа является взвешенной медианой с весами \(1/|y_i|\) и смещена в сторону меньших ответов. Отрицательные ответы требуют отдельного соглашения, которого у стандартного MAPE нет.

Симметричная относительная ошибка
\[
\mathrm{sMAPE}
=\frac1n\sum_{i=1}^{n}
\frac{2\lvert y_i-\hat y_i\rvert}{\lvert y_i\rvert+\lvert\hat y_i\rvert}
\]
содержит множитель \(2\) по соглашению; встречается запись без него. В протоколе указывают, какая именно формула вычислена.

\Ex При \(y=100\) и абсолютной ошибке \(50\) прогноз \(150\) даёт \(2\cdot 50/(100+150)=0{,}4\), прогноз \(50\) даёт \(2\cdot 50/(100+50)=2/3\). При одном и том же модуле ошибки больший прогноз увеличивает знаменатель и уменьшает штраф. Минимум по константе снова смещён вниз. Если \(y=\hat y=0\), знаменатель равен нулю; соответствующее слагаемое задают отдельным соглашением.

Среднеквадратическая логарифмическая ошибка
\[
\mathrm{RMSLE}
=\sqrt{\frac1n\sum_{i=1}^{n}
\bigl(\log(1+y_i)-\log(1+\hat y_i)\bigr)^{2}}
\]
сравнивает порядки величин. Единица под логарифмом допускает нулевое значение. Отрицательное значение не допускается.

\Ex При ответе \(100\) прогноз \(50\) даёт разность логарифмов около \(0{,}68\), прогноз \(150\) --- около \(0{,}40\). Абсолютная ошибка в обоих случаях равна \(50\), недопрогноз входит с большим весом. Если решение принимается в исходных единицах целевой переменной, RMSLE эту ошибку не выражает.

\subsection{Pinball-потеря}

Для \(\tau\in(0,1)\) и остатка \(r=y-\hat y\)
\[
\rho_{\tau}(r)=
\begin{cases}
\tau r, & r\ge 0,\\
(\tau-1)r, & r<0.
\end{cases}
\]
Положительный остаток означает, что прогноз меньше ответа.

\fig{figures/loss-pinball.png}{При \(\tau=0{,}8\) недопрогноз штрафуется сильнее перепрогноза.}

При \(\tau=1/2\) функция равна \(|r|/2\), и оптимальная константа есть медиана. При \(\tau=0{,}8\) одинаковые по модулю остатки неравноценны:
\[
\rho_{0{,}8}(2)=1{,}6,\qquad \rho_{0{,}8}(-2)=0{,}4.
\]
Оптимальная константа есть квантиль уровня \(\tau\). На выборке \(2,2,3,10\) большое \(\tau\) сдвигает её к крайнему наблюдению \(10\). Среднее \(4{,}25\) этот функционал не минимизирует. Тем самым pinball-потеря реализует третью характеристику из начала конспекта.

\subsection{Функция Хьюбера}

При пороге \(\delta>0\)
\[
H_{\delta}(r)=
\begin{cases}
\tfrac12 r^{2}, & |r|\le\delta,\\
\delta\bigl(|r|-\delta/2\bigr), & |r|>\delta.
\end{cases}
\]
В окрестности нуля функция квадратична и дифференцируема, за порогом растёт линейно. В точке \(|r|=\delta\) оба выражения совпадают.

\fig{figures/loss-huber.png}{При \(\delta=1\) центральная часть квадратична, хвосты линейны.}

\Ex При \(\delta=1\) и \(r=3\) продолжение квадратичной ветви дало бы \(\tfrac12\cdot 9=4{,}5\). Линейная ветвь даёт \(1\cdot(3-1/2)=2{,}5\). Удалённое наблюдение сохраняет влияние, но не квадратичное.

Функция Хьюбера используется как функция потерь при оценивании. В отчёт о качестве её не помещают: туда входят MAE или RMSE. Порог \(\delta\) выбирают по отложенной выборке и на контроле не пересматривают. Малые \(\delta\) приближают функционал к модулю, большие --- к квадрату.

\section{Логарифм целевой переменной}

Пусть оценивается модель \(\log y=x^{\top}w+\varepsilon\) при \(y>0\) и нормальном шуме \(\varepsilon\) с дисперсией \(\sigma^{2}\). Условное распределение \(y\) логнормально. Его медиана равна \(\exp(x^{\top}w)\), математическое ожидание ---
\begin{equation}
\label{eq:logmean}
\mathbb E[y\mid x]=\exp(x^{\top}w+\sigma^{2}/2).
\end{equation}
Слагаемое \(\sigma^{2}/2\) относится к нормальному шуму в логарифме. Универсальным переводным коэффициентом из логарифма в исходную шкалу оно не является.

\fig{figures/lognormal.png}{У логнормального распределения медиана расположена левее математического ожидания.}

Направление неравенства между средним и медианой следует из выпуклости экспоненты и без нормальности: при \(\mathbb E\varepsilon=0\) и непостоянном \(\varepsilon\) выполнено \(\mathbb E\exp(\varepsilon)>1\). Метрика, вычисленная по \(\log y\), и метрика, вычисленная по \(y\), отвечают на разные вопросы. Единица измерения фиксируется в протоколе.

\section{Сводка сквозного примера}

Промежуточные величины: \(\mathrm{RSS}=1/6\), \(\mathrm{SST}=2/3\), \(\bar y=5/3\), \(\hat w=(2/3,\ 1/2)\).

\begin{center}
\begin{tabular}{@{}lrr@{}}
\toprule
Метрика & Константа \(5/3\) & Прямая МНК\\
\midrule
MSE & \(2/9\) & \(1/18\)\\
RMSE & \(\sqrt{2/9}\) & \(\sqrt{1/18}\)\\
MAE & \(4/9\) & \(2/9\)\\
\(R^{2}\) & \(0\) & \(3/4\)\\
\(R^{2}_{\mathrm{adj}}\) & --- & \(1/2\)\\
MAPE & \(1/3\) & \(5/36\)\\
\bottomrule
\end{tabular}
\end{center}

Все строки вычислены на обучающих объектах. Добавление столбца не увеличивает \(\mathrm{RSS}\) на тех же объектах, а \(R^{2}\) строго возрастает, если новый столбец не ортогонален текущему остатку. Сравнение спецификаций поэтому переносят на объекты, не входившие в~\eqref{eq:normal}.

\section{Протокол оценивания}

Параметры \(w\), выборочное среднее константной модели, словарь категориальных уровней, а также среднее и стандартное отклонение любого преобразования признаков оцениваются на обучающей выборке.

На валидационной выборке выбирают состав столбцов, метрику отчёта и порог \(\delta\) функции Хьюбера. Контрольная выборка даёт одну оценку уже зафиксированного решения. Изменение спецификации после просмотра контроля лишает эту оценку смысла контроля: выборка становится ещё одной валидационной.

\Ex Следующие числа не продолжают сквозной пример. На валидации ответы равны \(6\) и \(14\), обучающее среднее равно \(10\). Модули константы равны \(4\) и \(4\), откуда \(\mathrm{MAE}=4\). Прогнозы прямой \(7\) и \(13\) дают модули \(1\) и \(1\), откуда \(\mathrm{MAE}=1\). Двух объектов недостаточно для заключения о пригодности модели. Пример показывает только устройство разреза.

\subsection{Схемы разбиения}

Разбиение должно воспроизводить условия будущего прогноза.

\fig{figures/splits.png}{Четыре схемы: случайное разбиение, сохранение долей корзин ответа, разбиение по времени и разбиение по группе.}

Случайное разбиение корректно, когда объекты независимы и одинаково распределены и не связаны общим временем или общей сущностью. Долю контроля фиксируют до вычисления метрики.

Если требуется сохранить доли дешёвых и дорогих объектов, непрерывный ответ заранее разбивают на корзины и в каждую часть помещают одинаковые доли корзин. Сырое значение \(y\) для стратификации непригодно: повторяющихся значений либо слишком много, либо нет вовсе.

При временной зависимости будущее не перемешивают с прошлым. Обучающая часть состоит из более ранних наблюдений, проверочная --- из более поздних. Перемешивание помещает будущий режим в оценку \(w\).

Объекты одной сущности --- сделки одного дома, кадры одного видео, строки одного пользователя --- не разделяют между обучением и проверкой. В противном случае алгоритм запоминает сущность, а разбиение этого не обнаруживает.

Обычная процедура \(k\)-fold, в которой блоки проверки чередуются, корректна в первой из четырёх ситуаций. Для времени используют разбиение, где обучение идёт только по прошлому, например \texttt{TimeSeriesSplit}. Для групп используют разбиение, где группа целиком попадает в одну часть, например \texttt{GroupKFold}.

\subsection{Утечки, существенные для этой лекции}

Утечкой называется присутствие в матрице плана сведения, которого не будет в момент применения алгоритма.

Если признак является функцией целевой переменной, например произведение цены квадратного метра на площадь, остаток на обучении может быть близок к нулю при отсутствии зависимости, доступной в момент прогноза.

Среднее для заполнения пропусков, параметры масштаба и словарь категорий оценивают на обучении и затем применяют к остальным частям. Пусть на обучении наблюдались \(0,0,10\), а отложенное значение равно \(10\). Обучающее среднее есть \(10/3\), среднее всех четырёх чисел есть \(5\). Одно и то же сырое значение вычитается из разных центров, и столбец, входящий в~\eqref{eq:normal}, оказывается другим. Второе среднее для преобразования уже непригодно.

Повтор одной сущности по обе стороны разбиения позволяет алгоритму запомнить сущность. Высокое значение метрики является поводом проверить список утечек. Само по себе оно не доказывает ни наличие утечки, ни качество модели.

\begin{warnbox}
Словарь one-hot, построенный по всей таблице, добавляет столбец уровня, которого в обучающей части могло не быть. Состав матрицы в нормальных уравнениях тогда другой. Уровень, встретившийся только в отложенной части, при словаре, оценённом на обучении, собственного коэффициента не имеет.
\end{warnbox}

\subsection{Запись эксперимента}

Одна строка соответствует одному законченному протоколу. Числа продолжают пример с ответами \(6\) и \(14\).

\begin{center}
\begin{tabular}{@{}p{48mm}p{34mm}p{28mm}r@{}}
\toprule
Гипотеза & Модель & Где оценена & MAE\\
\midrule
Достаточно константы & столбец единиц & обучение & 4\\
Линейная зависимость от площади & МНК, один признак & обучение & 1\\
Масштаб по всей таблице & тот же МНК & вся матрица & ---\\
\bottomrule
\end{tabular}
\end{center}

Первые две строки сравнимы: разбиение и метрика совпадают и выбраны заранее. Третья строка относится к другому протоколу и в сравнение не входит. В выводе указывают, улучшила ли прямая константную модель по заранее названной метрике, и в каких единицах выражено число.

\section{Что следует удерживать}

\begin{enumerate}
\item Восстанавливается условное среднее числового ответа. Простая и множественная модели различаются числом признаков; размерность ответа в обеих равна единице.
\item Параметры метода наименьших квадратов удовлетворяют системе \(X^{\top}X\,w=X^{\top}y\). При полном ранге решение единственно, при неполном единствен прогноз.
\item В сквозном примере \(w_0=2/3\), \(w_1=1/2\). Нулевая сумма остатков следует из столбца единиц.
\item Полные one-hot и one-cold вместе со свободным членом вырождены. Один уровень удаляют; его коэффициент читают как контраст с базовым уровнем.
\item Значение \(R^{2}=3/4\) в сквозном примере есть доля суммы квадратов на обучающих объектах.
\item Функция потерь и метрика отчёта могут различаться. Протокол называет обе.
\item Преобразования признаков и словарь уровней оцениваются на обучении. Контроль после выбора спецификации не используется для её пересмотра.
\end{enumerate}

\section{Контрольные вопросы}

\begin{enumerate}
\item Почему полный one-hot из трёх уровней вместе со свободным членом не задаёт единственного \(w\)?
\item Какой столбец матрицы плана обеспечивает равенство нулю суммы остатков в сквозном примере?
\item Почему \(R^{2}=3/4\) на этих трёх объектах не является оценкой качества?
\item Является ли ошибкой протокола оценивание по квадрату при отчёте по MAE?
\item Что происходит с временным рядом цен, если объекты перемешать перед разбиением?
\item Чему равно \(\exp(x^{\top}w)\) после оценивания модели для \(\log y\), и при каком предположении о шуме известно условное среднее?
\item В каком смысле полиномиальная модель остаётся линейной, хотя график \(a(x)\) криволинеен?
\item Совпадает ли задача логистической регрессии с задачей~\eqref{eq:cond-mean}?
\item Если два столбца \(X\) совпали, теряет ли единственность вектор прогнозов на обучении?
\item Какая константа минимизирует средний квадрат, и из какого частного случая нормальных уравнений она получается?
\end{enumerate}

\subsection{Ответы}

\begin{enumerate}
\item Сумма трёх индикаторов совпадает со столбцом единиц, поэтому ранг \(X\) меньше \(p\). Прибавление вектора из ядра меняет \(w\) и не меняет \(Xw\).
\item Столбец единиц: остаток ему ортогонален. В числах \(-1/6+1/3-1/6=0\).
\item Коэффициент вычислен на объектах, по которым искали \(w\). На новых объектах значение может быть иным, в том числе отрицательным. Кроме того, \(R^{2}\) есть доля суммы квадратов, а не доля точных ответов.
\item Нет, если протокол раздельно указывает функцию потерь и метрику. Минимум MAE метод наименьших квадратов не обеспечивает. Совпадение направления сравнения в сквозном примере, \(2/9\) против \(4/9\), есть свойство этой выборки.
\item Наблюдения из будущего попадают в оценку параметров. Для временного ряда обучающая часть предшествует проверочной.
\item При нормальном шуме в логарифме это условная медиана \(y\). Условное среднее даётся формулой~\eqref{eq:logmean}.
\item Столбцы \(x\) и \(x^{2}\) фиксируются до оценивания. По параметру \(w\) алгоритм остаётся скалярным произведением \(x^{\top}w\).
\item Нет. Логистическая регрессия относится к классификации и оценивает вероятность метки.
\item Нет. Совпадение столбцов уничтожает единственность \(w\) и интерпретацию отдельного коэффициента. Прогноз есть проекция на оболочку столбцов и прибавлением вектора из ядра не меняется.
\item Выборочное среднее обучающих ответов. В модели из столбца единиц условие нулевой суммы остатков даёт \(w_0=\bar y\).
\end{enumerate}

\section{Задача}

Для сквозного примера и константы \(\hat y=5/3\) выполните следующее.

\begin{enumerate}
\item Вычислите MAE и MSE обеих моделей и сверьте результат со сводной таблицей.
\item Добавьте объект \((x,y)=(4,10)\) только в контроль. Параметры прямой и константу не переоценивайте.
\item Сравните на этом объекте абсолютные ошибки прямой и константы.
\item Укажите, какое разбиение здесь реализовано, и почему одного контрольного объекта недостаточно для вывода.
\end{enumerate}

\subsection{Решение}

Первый пункт совпадает со сводкой: у константы \(\mathrm{MAE}=4/9\) и \(\mathrm{MSE}=2/9\), у прямой \(\mathrm{MAE}=2/9\) и \(\mathrm{MSE}=1/18\).

Параметры не переоцениваются. При \(x=4\)
\[
\hat y=\frac23+\frac12\cdot 4=\frac83.
\]
Абсолютные ошибки равны
\[
\Bigl|10-\frac83\Bigr|=\frac{22}{3},
\qquad
\Bigl|10-\frac53\Bigr|=\frac{25}{3}.
\]
На данном объекте прямая точнее. Объект не участвовал в системе~\eqref{eq:normal}, поэтому сравнение устроено корректно. Одного объекта недостаточно: следующий контрольный объект может изменить знак разности ошибок. В протоколе такой результат описывают как сравнение на одном наблюдении, без заключения о качестве модели в целом.

\end{document}
